\begin{theorem}[Jensen’s inequality]
Let $f:\mathbb R\to\mathbb R$ be convex.  
For any integer $n\ge 1$, any real numbers $x_{1},\dots ,x_{n}$ and any
weights $\lambda_{1},\dots ,\lambda_{n}$ satisfying
\[
\lambda_i\ge 0\;(i=1,\dots ,n),\qquad
\sum_{i=1}^{n}\lambda_i =1,
\]
we have
\begin{equation}
f\!\Bigl(\sum_{i=1}^{n}\lambda_i x_i\Bigr)\le
\sum_{i=1}^{n}\lambda_i f(x_i).\tag{1}
\end{equation}
Equality holds exactly when either all points with positive weight are
identical, or $f$ is affine on the convex hull of those points.
\end{theorem}

\begin{proof}
\textbf{Binary convexity.}  
Convexity of $f$ is equivalent to
\begin{equation}
f\bigl(t x+(1-t)y\bigr)\le t\,f(x)+(1-t)f(y)\qquad
(x,y\in\mathbb R,\;t\in[0,1]).\tag{C}
\end{equation}

\medskip\noindent\textbf{Base case $n=1$.}  
If $n=1$, then the only admissible weight is $\lambda_{1}=1$, so
\[
f\!\Bigl(\sum_{i=1}^{1}\lambda_i x_i\Bigr)=f(x_{1})
=\sum_{i=1}^{1}\lambda_i f(x_i),
\]
and (1) holds with equality.

\medskip\noindent\textbf{Induction hypothesis.}  
Assume that for some $k\ge1$ inequality (1) holds for every choice of
$x_{1},\dots ,x_{k}$ and non‑negative weights $\lambda_{1},\dots ,\lambda_{k}$
with $\sum_{i=1}^{k}\lambda_i=1$:
\begin{equation}
f\!\Bigl(\sum_{i=1}^{k}\lambda_i x_i\Bigr)\le
\sum_{i=1}^{k}\lambda_i f(x_i).\tag{IH}
\end{equation}

\medskip\noindent\textbf{Inductive step $k\to k+1$.}  
Let $\lambda_{1},\dots ,\lambda_{k+1}\ge0$ with $\sum_{i=1}^{k+1}\lambda_i=1$
and let $x_{1},\dots ,x_{k+1}\in\mathbb R$.

Set
\[
\mu:=\sum_{i=1}^{k}\lambda_i,\qquad
\lambda_{k+1}=1-\mu .
\]

If $\mu=0$ then $\lambda_{k+1}=1$ and we are back to the base case,
so (1) is true.  Henceforth assume $\mu>0$.

Define normalized weights for the first $k$ points,
\[
\alpha_i:=\frac{\lambda_i}{\mu}\qquad(i=1,\dots ,k),
\]
and the intermediate convex combination
\[
y:=\sum_{i=1}^{k}\alpha_i x_i
   =\frac{1}{\mu}\sum_{i=1}^{k}\lambda_i x_i .
\]
Clearly $\alpha_i\ge0$ and $\sum_{i=1}^{k}\alpha_i=1$.

\medskip\noindent\textbf{Step 1 – binary convexity.}  
Apply \eqref{C} with $x:=y$, $y:=x_{k+1}$ and $t:=\mu\in[0,1]$:
\begin{equation}
f\bigl(\mu y+(1-\mu)x_{k+1}\bigr)
\le \mu f(y)+(1-\mu)f(x_{k+1}).\tag{2}
\end{equation}
Since $\mu y+(1-\mu)x_{k+1}
      =\sum_{i=1}^{k}\lambda_i x_i+\lambda_{k+1}x_{k+1}
      =\sum_{i=1}^{k+1}\lambda_i x_i$,
inequality \eqref{2} becomes
\begin{equation}
f\!\Bigl(\sum_{i=1}^{k+1}\lambda_i x_i\Bigr)
\le \mu f(y)+(1-\mu)f(x_{k+1}).\tag{3}
\end{equation}

\medskip\noindent\textbf{Step 2 – induction hypothesis on the first $k$ points.}  
Applying (IH) to $x_{1},\dots ,x_{k}$ with weights $\alpha_{1},\dots ,\alpha_{k}$ yields
\[
f(y)=f\!\Bigl(\sum_{i=1}^{k}\alpha_i x_i\Bigr)
\le\sum_{i=1}^{k}\alpha_i f(x_i).\tag{4}
\]
Multiplying (4) by $\mu$ and using $\mu\alpha_i=\lambda_i$ gives
\begin{equation}
\mu f(y)\le\sum_{i=1}^{k}\lambda_i f(x_i).\tag{5}
\end{equation}

\medskip\noindent\textbf{Step 3 – conclusion of the inductive step.}  
Insert (5) into (3):
\[
f\!\Bigl(\sum_{i=1}^{k+1}\lambda_i x_i\Bigr)
\le\sum_{i=1}^{k}\lambda_i f(x_i)+(1-\mu)f(x_{k+1})
 =\sum_{i=1}^{k+1}\lambda_i f(x_i),
\]
which is precisely (1) for $n=k+1$.

Thus the inequality holds for all $n\ge1$ by induction.

\medskip\noindent\textbf{Equality cases.}  
In the binary convexity \eqref{C} equality occurs iff either $t=0$ or $t=1$
(one of the two weights is zero) or $f$ is affine on the segment joining the
two points.  Equality in the induction hypothesis (IH) occurs iff either
all points with positive weight among $x_{1},\dots ,x_{k}$ coincide or
$f$ is affine on their convex hull.

Combining these observations through the inductive construction shows
that equality in \eqref{1} holds exactly when

\begin{enumerate}[(i)]
\item all points $x_i$ with $\lambda_i>0$ are identical, or
\item $f$ is affine on the convex hull $\operatorname{conv}\{x_i:\lambda_i>0\}$.
\end{enumerate}
These are the classical equality conditions for Jensen’s inequality.
\qed
\end{proof}